\section{Literature Review: From Externalities to Incentives}

Lorem ipsum dolor sit amet, consectetur adipiscing elit. Sed do eiusmod tempor incididunt ut labore et dolore magna aliqua. Ut enim ad minim veniam, quis nostrud exercitation ullamco laboris nisi ut aliquip ex ea commodo consequat. Duis aute irure dolor in reprehenderit in voluptate velit esse cillum dolore eu fugiat nulla pariatur. Excepteur sint occaecat cupidatat non proident, sunt in culpa qui officia deserunt mollit anim id est laborum.

% A related difficulty emerges in the way production theory represents nature. A traditional strand of economic analysis defines the factors of production as land or natural resources (N), capital (K), and labor (L), the fundamental inputs in the production of goods and services. Although ecological conditions are directly or indirectly embedded in all three, dominant strands of economic modeling have often treated the environment either as an exogenous constraint, an externality, a fixed natural endowment, or a corrective term, rather than as an endogenous and dynamic component of production itself \cite{jehle2011}. As a result, ecological constraints are frequently ignored or only indirectly considered through land productivity, which affects land rents according to natural attributes such as fertility, location, and regenerative capacity \cite{hubacek2002role}.

% A related difficulty emerges in the way production theory represents nature. A traditional strand of economic analysis defines the factors of production as land or natural resources (N), capital (K), and labor (L), the fundamental inputs in the production of goods and services. Although ecological conditions are directly or indirectly embedded in all three, dominant strands of economic modeling have often treated the environment either as an exogenous constraint, an externality, a fixed natural endowment, or a corrective term, rather than as an endogenous and dynamic component of production itself \cite{jehle2011}. As a result, ecological constraints are frequently ignored or only indirectly considered through land productivity, which affects land rents according to natural attributes such as fertility, location, and regenerative capacity \cite{hubacek2002role}.

% A growing body of literature has sought to address this gap by endogenizing nature within production theory, recognizing that ecological conditions can directly influence the productivity of capital and labor, and that economic activities can, in turn, affect ecological dynamics. This approach emphasizes the need for economic models to incorporate feedback loops between economic decisions and ecological outcomes, moving beyond the traditional externality framework to a more integrated understanding of socio-ecological systems \cite{costanza1997value, daily1997nature, farley2012ecosystem}.

% Alternatively, \citeonline{ostrom2002commonpool} challenged Hardin’s dilemma by arguing that, under certain institutional conditions, the self-management of natural resources can be sustainable. \citeonline{assuncao2013deterring} highlight that enforcement, when paired with well-designed incentives, has shown results in reducing Amazon deforestation, yet faces structural and political resistance when scaled. \citeonline{scheinkman2008social}’s contributions to the modeling of local interactions and externalities provide a relevant foundation for understanding how decentralized behavior can lead to socially suboptimal equilibria, especially in resource use contexts. 
